\section{Introduction}
\label{sec:introduction}

\subsection{Motivation and Context}
\label{sec:intro_motivation}

Indoor air pollution contributes to approximately 3.2 million premature deaths annually, making continuous monitoring of pollutants such as PM$_{2.5}$, PM$_{10}$, volatile organic compounds (VOCs), CO$_2$, CO, and NO$_2$ a pressing public health imperative~\cite{karmakar2024exploring}. The proliferation of low-cost IoT sensor platforms has enabled fine-grained spatio-temporal mapping of indoor air quality (IAQ) in real time~\cite{dimitriou2025innovations, enviroiot2025}. However, these systems face a fundamental tension between monitoring fidelity and energy sustainability~\cite{energysustainable2023}. Higher sampling rates and greater ADC resolution improve pollutant characterisation accuracy but increase power draw at the sensor, data volume for transmission, and computational load for on-device inference.

Edge devices deployed for air quality monitoring---typically built around microcontrollers such as the ESP32---operate under stringent resource constraints. The ESP32 consumes approximately 45.5~mA (0.2275~W) in active mode and 5.27~mA (0.0264~W) in sleep mode at 5~V~\cite{esp32power2025}. With deep sleep, consumption drops to 0.05~mA, enabling battery lifetimes exceeding 48 months under intermittent duty cycling~\cite{benchmarking2022}. However, continuous 1~Hz multi-channel monitoring demands sustained active-mode operation, rapidly depleting energy budgets. The ADC contributes significantly to this overhead; its power scales exponentially with bit resolution according to the Walden figure of merit (FoM), where state-of-the-art successive-approximation-register (SAR) ADCs achieve 0.44--16.6~fJ per conversion step~\cite{waldenADC2018, andrulis2024adc}.

\subsection{Problem Statement}
\label{sec:intro_problem}

Neural network quantisation has emerged as a powerful technique for reducing inference energy on edge devices. Post-training quantisation (PTQ) compresses weights and activations from 32-bit floating point to lower precision---typically INT8---without expensive retraining~\cite{nagel2021whitepaper, quantsurvey2021}. SmoothQuant~(2024) demonstrated training-free W8A8 quantisation that preserves accuracy through mathematically equivalent transformations~\cite{smoothquant2024}. INT8 PTQ has achieved 14.5$\times$ speedup at only 1.96~W average power on edge hardware~\cite{palakodety2025tpu}. At 4-bit precision, fixed-point designs maintain 99.28\% accuracy with less than 0.3\% degradation~\cite{guo2020fixedpoint}.

Despite these advances, sensor-level quantisation and inference-level quantisation have been treated as independent problems. Sensor data compression methods---adaptive thresholding~\cite{atitallah2023autonomous}, bounded-error pruning~\cite{azar2020bounded}, and SAX-based reduction~\cite{aljanabi2022sax}---operate on raw signals without considering downstream inference. Model compression techniques such as PTQ~\cite{smoothquant2024}, quantisation-aware training (QAT)~\cite{edgeqat2024, efficientqat2024}, and mixed-precision quantisation~\cite{dong2024efficient, dybit2023} optimise the neural network in isolation from the sensing pipeline. This separation leaves significant energy savings unrealised. Q-learning-based adaptive data delivery~\cite{qlearning2025airquality} reduces transmission by 40--50\% but operates solely at the communication layer. Resolution reconfiguration guided by machine learning has shown 8.7--23$\times$ power savings in neural recording~\cite{pandey2025neural}, demonstrating the potential of per-channel precision control, yet this concept has not been applied to environmental IoT monitoring.

\subsection{Contributions}
\label{sec:intro_contributions}

This paper makes four contributions:
\begin{enumerate}[leftmargin=*]
    \item A Variance-Aware Adaptive Bit-Width Quantisation (VABQ) framework that unifies sensor-level ADC resolution control with inference-level PTQ in a single optimised pipeline.
    \item A variance-entropy criterion that dynamically selects between 4-bit, 8-bit, and 16-bit ADC resolution based on rolling signal statistics, exploiting the periodic structure of indoor pollutant data.
    \item Comprehensive experiments on the DALTON dataset~\cite{karmakar2024exploring, karmakar2024indoor} demonstrating 61.3\% sensor-level and 74.8\% inference-level energy reductions with less than 3.9\% accuracy degradation.
    \item A comparative analysis against four baseline methods---uniform 16-bit, static 8-bit PTQ, fixed-threshold adaptive sampling, and SAX-LZW compression---validating the superiority of the dual-layer approach.
\end{enumerate}

\subsection{Paper Organisation}
\label{sec:intro_organisation}

The remainder of this paper is organised as follows. Section~\ref{sec:related_work} reviews related work and identifies the research gap. Section~\ref{sec:methodology} details the VABQ methodology. Section~\ref{sec:experiments} describes the experimental setup and dataset. Section~\ref{sec:results} presents quantitative results. Section~\ref{sec:discussion} discusses practical implications, novelty justification, and limitations. Section~\ref{sec:conclusion} concludes the paper.
